% ============ 07. Plan ============
\section{Plan de acción para llegar a ``listo'' (3--5 días dirigidos)}

Índice propuesto del anteproyecto (consenso de los evaluadores), con la fuente
de cada sección:

\begin{enumerate}
  \item \textbf{Problema y contexto — 1½--2 págs. (D1).} Redactar actor, alcance,
        frecuencia, consecuencia y pregunta estadística, usando cifras propias
        verificadas (51\,\% Bogotá+Antioquia; 24\,\% vs.\ 48\,\% mujeres por
        área; 3--8$\times$ vs.\ DANE) y una fuente externa de magnitud en
        Colombia. \emph{Fuente: informe base \S1 + datos de auditoría.}
  \item \textbf{Estado del arte — 2--3 págs. (D2).} Sintetizar 6--7 enfoques del
        SotA (20 págs. $\rightarrow$ 3), añadir contraste explícito de 2
        aproximaciones y el vacío declarado; convertir a APA 7.
  \item \textbf{Objetivos — ½ pág. (D3).} Reformular OG + 3--5 OE como logros
        verificables alineados al problema. \emph{Fuente: informe base \S3
        (reescribir).}
  \item \textbf{Datos y método — 1--2 págs. (D4).} Socrata
        \texttt{bqtm-4y2h}/\texttt{33dq-ab5a} (procedencia, periodo, unidad,
        volumen, acceso y \emph{licencia}), tabla objetivo$\rightarrow$
        variables$\rightarrow$ método, una alternativa descartada (p.~ej.,
        encuesta primaria o PySpark), plan B si el dato no llega (redescarga
        API/DVC) y novedad frente a lo publicado.
  \item \textbf{Alcance, riesgos y cronograma — 1 pág. (D5).} Exclusiones
        (informe base \S4), 12 semanas en hitos con entregables, 2 de holgura,
        2--3 riesgos con contingencia y horas por integrante.
  \item \textbf{Declaración de IA y reparto — ½ pág. (S3).} Usar el borrador del
        agente S, firmarlo con nombres/roles y re-verificar los casos.
  \item \textbf{Referencias — APA 7} (no cuenta en el total; ~18--25 citadas).
\end{enumerate}

\textbf{Preparación oral (S1/S2):} banco de preguntas con respuestas propias
(10 ya propuestas), guion de 10 minutos y ~9 diapositivas; práctica individual
de que cada integrante defienda el conjunto; verificación de que la declaración
de IA es creíble frente a un interlocutor escéptico.

\begin{tcolorbox}[hallazgo]
\textbf{Conclusión final:} los materiales actuales valen \emph{más que un
borrador} — son la evidencia verificada que la rúbrica premia (cifras propias,
0 referencias fabricadas, datos reales abiertos y verificados). El trabajo que
queda es de \textbf{redacción dirigida y decisión}, no de investigación: fijar
el objeto, escribir problema y objetivos, sintetizar el SotA a APA 7, producir
la declaración de IA con nombres y practicar la sustentación individual.
Cumpliendo ese plan, el anteproyecto apunta a \textbf{Competente--Ejemplar} en
los 8 criterios, sin riesgo en los bloqueantes.
\end{tcolorbox}
